\section{Prospective Evaluation}

The primary comparison is M2 against the rolling 30-day rate, while M2 against global M0 is retained
for compatibility. The theory-feature test asks whether M3-lite improves over M2. Log loss is the key secondary metric; 168-hour scores, calibration,
and PR-AUC are secondary analyses. Uncertainty uses paired seven-day block bootstrap with 2,000
resamples.

The first formal comparison requires both 180 valid scheduled days and 20 prospective 24-hour
positives. A historical block-variance calculation suggests that reliably detecting a Brier
difference of 0.004 may require substantially more than 180 days, so this minimum is not presented
as a power guarantee. Bootstrap runs, excluded runs, and missed runs do
not count toward the primary sequence. At the time of this draft, no valid scheduled run has matured.
The prospective result is therefore intentionally blank. Until the stopping rule is reached, the
dashboard reports operational completeness and descriptive cumulative scores without declaring a
winner.
